% answers_short.tex —— 2015 年·福建卷·文科 · 简短答案速查
% 由 tools/make_short_answers.py 自 answers.tex 的 \gksolhead/\gkanswer 提取，
% 并与 2015/2015_数学真题_福建文_答案要点.md 逐题比对后生成，勿手改（改源 MD 后重跑）。

\begin{gknotice}
\textbf{使用说明}：本节只给最终答案便于快速核对，完整推导见解析版。答案为编写者独立推导，\textbf{非官方标准答案}。
\end{gknotice}

\gkpart{选择题}
\begin{center}\begin{tabular}{ccccccccccccc}
\toprule
题号 & 1 & 2 & 3 & 4 & 5 & 6 & 7 & 8 & 9 & 10 & 11 & 12 \\
\midrule
答案 & A & D & D & C & C & D & A & B & B & C & A & B \\
\bottomrule
\end{tabular}\end{center}

\gkpart{填空题}
\begin{itemize}[leftmargin=2.6em]
  \item[\textbf{13.}] \(25\)
  \item[\textbf{14.}] \(\sqrt2\)
  \item[\textbf{15.}] \(1\)
  \item[\textbf{16.}] \(9\)
\end{itemize}

\gkpart{解答题}
\begin{itemize}[leftmargin=2.6em]
  \item[\textbf{17.}] (1) \(a_n=n+2\)；(2) \(2101\)
  \item[\textbf{18.}] (1) \(\dfrac9{10}\)；(2) \(\dfrac{121}{20}=6.05\)
  \item[\textbf{19.}] (1) \(y^2=4x\)；(2) \(d(F,GA)=d(F,GB)=\dfrac{4\sqrt2}{\sqrt{17}}\)，故两直线同为以 \(F\) 为圆心圆的切线
  \item[\textbf{20.}] (1) 结论成立；(2) 最大体积为 \(\dfrac13\)；(3) 最小值为 \(\dfrac{1+\sqrt3}{\sqrt2}\)
  \item[\textbf{21.}] (1) 最小正周期 \(2\pi\)；(2) ① \(g(x)=10\sin x-8\)。② 存在无穷多个互不相同的正整数 \(x_0\)，使 \(g(x_0)>0\)
  \item[\textbf{22.}] (1) 单调递增区间为 \(\left(0,\dfrac{1+\sqrt5}{2}\right]\)；(2) 结论成立：当 \(x>1\) 时，\(f(x)<x-1\)；(3) \(k<1\)
\end{itemize}
